\section{A separate Grothendieck dictionary: objects, operations, and proof transport}
\label{sec:grothendieck-dictionary}

This part restates the checked deductions in the requested geometric language. A locator identifies the framework in which a proof lives; it is not a claim that Grothendieck stated a modern complexity result. The original manuscript images are not in the attachment. Cote and scan references below therefore refer to the supplied transcriptions, and the mathematical corrections do not certify an unseen handwritten reading.

\subsection{Exact framework locators and scope}

\paragraph{D-S3: group schemes and weights.}
For an affine group scheme, its representations are comodules over its Hopf algebra. The relevant reference is SGA~3, Expos\'e~I, Proposition~4.7.2. For a diagonalizable group $D(M)$, such comodules are $M$-graded modules, by Proposition~4.7.3; the following remark gives the corresponding graded-algebra description of affine schemes with action. Proposition~4.7.4 gives equivariant splittings when the underlying sequence splits. These statements, inspected in the original Expos\'e~I PDF, give the integral meaning of every torus weight and cocharacter below. For stabilizers, SGA~3, Expos\'e~VI$_B$, Proposition~4.1 gives upper semicontinuity of group-fiber dimension and explicitly cites EGA~IV$_3$, Theorem~13.1.3. The latter locator is corroborated through that primary SGA text; its entire EGA proof has not been reread here.

\paragraph{D-S6: $K_0$, exterior powers, Adams operations, and filtration.}
SGA~6, Expos\'e~V, Definition~2.1 and Example~2.2(b), construct the pre-$\lambda$ structure of locally projective modules over a ringed topos; Lemma~2.2.1 uses a filtration on exterior powers of an exact sequence. Definition~2.4 and equation~(2.4.2) specify the universal product and composition identities. Expos\'e~VI, Corollary~3.2 is the special-$\lambda$-ring assertion for the geometric $K'$ in its stated setting. Expos\'e~V, Proposition~3.2 defines $\gamma$ from $\lambda$; section~3.10, equation~(3.10.1), defines the augmentation filtration. Section~7.1, equation~(7.1.1), and Proposition~7.5 give Adams operations and their ring identities. The electronic transcriptions of the cited Expos\'e~V passages were inspected; the exact VI~3.2 locator was located in the text's index and corroborating primary references. We do not describe an integral character identity as an integral semisimple decomposition.

For split reductive group schemes a useful complementary reference is Serre (1968), Theorem~5, section~3.7, and the explicit $\mathrm{GL}_n$ character ring in section~3.8. The base-change statements in that theorem concern the representation Grothendieck group, not preservation of a characteristic-zero basis of simple objects in every characteristic. Nor do they automatically answer the separate finite-symmetric-group lattice question recorded in the supplied study of cote~133.

\paragraph{D-FGA: images, quotients, and projective parameter spaces.}
For a homogeneous orbit map over $\mathbb Z$, take its schematic image by the kernel of the coordinate homomorphism, as in G-04. Degreewise kernels and images are given by finite matrices and their minors. EGA~II, sections~2.4--2.5, in particular Proposition~2.5.1, gives the projective scheme and associated sheaf of a graded module. The electronic EGA translation was inspected for this locator. FGA, Bourbaki Expos\'e~221, Theorem~3.1 gives the Quot parameter space, Theorem~3.2 its fixed-Hilbert-polynomial projective form, and Proposition~3.8 its Grassmannian realization in the required degree range. The original FGA scan was inspected. These results place the families in parameter spaces; they supply no polynomial bound for finding the relevant minors or their nonvanishing loci.

\paragraph{D-ET: traces, weights, and their arithmetic size.}
SGA~4$\frac12$, \emph{Rapport sur la formule des traces}, Theorems~3.1--3.2 and Proposition~3.3, gives the determinant expression for the $L$-function, the trace formula, and the logarithmic-derivative identity. SGA~5, Expos\'e~X, Theorem~7.1 and Corollary~7.12, formula~(7.13), gives the Euler--Poincar\'e conductor accounting. Deligne's Weil~II, Theorem~3.3.1 and section~3.3.10, supplies the compact-support weight estimate used in P-06; Theorem~3.5.3 gives the equidistribution setting. The Weil~II locators were checked against Katz and Sarnak (1999), sections~9.0.13--9.0.16 and Theorem~9.2.6. The original large Weil~II PDF was not successfully opened here. This is a statement-and-dependency check, not an independent verification of the complete proof of Weil~II.

\subsection{Rectangular characters and orbit-closure multiplicities}

\paragraph{R-01. Tensor pullback of a Schur operation.}
For locally free $E,F$ of ranks $a,n$, the tensor morphism $\mathrm{GL}(E)\times\mathrm{GL}(F)\to\mathrm{GL}(E\otimes F)$ is defined over $\mathbb Z$. The class $s_{(d^n)}[E\otimes F]$ restricts, after passage to characteristic zero, to the sum of $[S_\mu E\otimes S_\nu F]$ with coefficients $g((d^n),\mu,\nu)$. The Haar projector is the complex realization of extraction of that character component. Its finite CUE dimension is a rank of a Schur construction; its normalized limiting constant need not be an integer rank. This is precisely the proof of R-01 transported through D-S3 and D-S6. Cote~133, scan~028, lines~6--17, provides the tensor-power and symmetric-operation language, while scan~019 distinguishes rank from the full character class.

\paragraph{R-02. Two different meanings of positivity.}
The tableau expansion is an integral character identity in the universal special $\lambda$-ring, with the Weyl alternating factor extracting a highest weight. Over $\mathbb C$, the invariant summand of $S_\lambda(E\otimes F)\otimes\det(E)^{-d}\otimes\det(F)^{-d}$ is an actual vector space. A Hermitian metric identifies its projector trace with a sum of squared absolute values. The former is an integral identity; the latter uses complex conjugation and compact integration. D-S6 and cote~133, scan~028, lines~19--24 (the direct-summand question), are the exact conceptual distinction. No cancellation-free combinatorial basis is inferred from the analytic squares.

\paragraph{R-03. Central quotient, disconnected stabilizer, and rank.}
The tensor morphism has the central subgroup $t\mapsto(tI,t^{-1}I)$ in its kernel, a statement of group schemes over $\mathbb Z$. A compact block-Levi branching rule is therefore a different restriction functor. The transpose component of the determinant stabilizer is additional group-scheme data beyond the connected tensor subgroup. The $S_4$ counterexamples compare distinct invariant functors and survive as finite character computations. D-S3 places these groups and actions; cote~133, scan~019, lines~45--56, is the corresponding warning that augmentation loses equivariant information. The negative Weingarten coefficient is a complex integration fact, not a negative rank of a locally free sheaf.

\paragraph{R-04. Plethysm as an operation on representation $K_0$.}
The cycle-index formula expresses $\sigma^d(\sigma^n[E])$ using the Adams operations $\psi^j$. D-S6, V~7.1 and V~7.5, makes their algebraic status precise. They are ring operations; $g\mapsto g^j$ need not be a group homomorphism, and $\psi^j$ need not preserve the cone of honest representations. The $U(2)$ mixed-moment counterexample therefore translates to the statement that separate character distributions do not determine the joint operation. The wreath-subgroup invariants give the characteristic-zero tensor-functor realization. Cote~133, scan~028, lines~12--17 and~26--35, places tensor powers and induction products, with each $\rho_i$ in its own homogeneous degree.

\paragraph{R-05. Image modules, not merely stabilizer invariants.}
The intrinsic object is the image of the degree-$d$ orbit-coefficient map. Its dual is the span of the orbit of $f^{\otimes d}$. Over $\mathbb C$, complete reducibility identifies the multiplicity in each Schur summand with the rank of its multiplicity matrix. The Haar covariance and its reduced Gram matrices have exactly those images; their principal exterior coefficients are squared Pl\"ucker norms. This transports the entire rank proof to determinantal rank strata and Grassmannians in D-FGA. The Hermitian form is an auxiliary real structure and is not asserted to define a motive. Cote~133, scan~028's question about a direct summand is answered here only after passing to characteristic zero.

\paragraph{R-06. Integral matrices and special fibers.}
The polynomial orbit map is defined over $\mathbb Z$ after the specified symmetrizer scalings. Its finite coefficient matrix is an integral map of free modules. The Gaussian covariance is another complex Gram form on its image, while the integral character-sum matrix computes the required Schur-isotypic rank after scalar extension. The Schur projectors require the stated factorial denominators. G-04 specifies the open localization on which the ranks and kernels specialize correctly. D-FGA records the associated rank and Quot strata; D-S6 records the tensor and Schur operations. This translation preserves the exponential matrix sizes, so changing language does not remove the computational gate.

\paragraph{R-07. Binary-form restriction and highest-weight sections.}
The determinant-pencil substitution is a morphism of coefficient schemes, defined after clearing the displayed binomial denominators. The sections $Q_j$ are highest-weight polynomial functions, and independence is proved by their restrictions to the chosen diagonal-pencil chart. Over $\mathbb C$ every binary form is a product of linear forms, so the binary coefficient subscheme lies in the determinant orbit closure. Highest-weight sections of length at most two therefore cannot belong to its ideal. This is the same proof expressed as ideal restriction in D-FGA and weights in D-S3. Cote~133, scan~028 supplies the polynomial-functor viewpoint; it does not supply a new multiplicity formula by rank alone.

\paragraph{R-08. The quotient class in equivariant $K_0$.}
For the polynomial ring $S=\operatorname{Sym}(W)$ and orbit quotient $A=S/I$, an equivariant free resolution yields $[A]=\sum_i(-1)^i[S\otimes T_i]$ in the appropriate graded Grothendieck group. Symmetric powers produce the Molien denominator and the alternating resolution produces the numerator. Exactness, not pointwise sign, proves the formula. D-S6, V~2.2(b), and D-FGA place the additivity and coherent quotient. Cote~133, scan~019, lines~9--20, contains the exterior/symmetric generating-series relation. Ignoring the numerator forgets the defining ideal and is not an operation allowed by $K_0$ additivity.

\paragraph{R-09. A four-factor invariant and a finite-rank cutoff.}
The squared rectangular multiplicity sum is the determinant-isotypic component of $\operatorname{Sym}(E_1\otimes E_2\otimes E_3\otimes E_4)$ in degree $nd$, obtained by composing two tensor restriction functors. The $U(N)$ cutoff means that Schur functors with more than $N$ rows vanish on a rank-$N$ bundle. The stable ratios denominator is therefore an identity of symmetric operations on four tensor factors. D-S3 and D-S6 carry the proof without an arithmetic purity hypothesis. Cote~133, scan~028, is the corresponding tensor-power locus. Nonnegative character norms are exact, but no effective algorithm for the remaining coefficient extraction has been added by the translation.

\paragraph{R-10. The quadratic case and the component group.}
Over $\mathbb C$, the $2\times2$ determinant is a nondegenerate quadratic form in four variables. Its orbit closure is the whole quadratic-form coefficient space. The connected group is the tensor image of $\mathrm{SL}_2\times\mathrm{SL}_2$; the full quadratic stabilizer includes the other orthogonal component. Accordingly the all-odd four-row part present in the ordinary rectangular Kronecker coefficient is removed from the symmetric coefficient and the coordinate ring. D-S3 records the group and component distinction, while D-S6 identifies the even-part Schur decomposition. Cote~133, scan~019's augmentation does not retain this distinction.

\paragraph{R-11. Low-degree surfaces and restricted ideals.}
The linear and quadratic determinant constructions are coefficient-scheme morphisms. The cubic statement uses a dense open set of smooth cubic surfaces and Beauville's determinantal representation theorem; closedness extends the containment. Block padding is a morphism over the coefficient ring, and the highest-weight ideal inclusion is the same D-S3/D-FGA argument as G-08. It excludes four-row multiplicity obstructions for padded forms of degree at most three without asserting a global orbit containment for arbitrary many-variable cubics. Cote~133, scan~028 provides the Schur-functor interpretation of the resulting coordinate modules.

\paragraph{R-12. A Jacobian certificate and its characteristic-zero consequence.}
The permanent-pencil coefficient map is a morphism of affine schemes over $\mathbb Z$. A maximal Jacobian minor nonzero modulo the stated prime is a nonzero integer and therefore nonzero over $\mathbb Q$. In characteristic zero the differential-rank criterion gives dominance, so the four-variable projection of the permanent orbit closure fills the coefficient space in the certified sizes. D-FGA then translates the hypersurface quotient sequence of G-07 into the exact four-row multiplicity difference. The finite-field calculation is a certificate for an integer minor; it is not an appeal to equidistribution or an assertion of semisimplicity modulo that prime. The tensor-power language is cote~133, scan~028; the rank/character distinction remains scan~019.

\subsection{Isolation as a torus action and a problem of coefficient modules}

\paragraph{I-01. Independent variable coordinates.}
The vertex-variable path polynomial is a section of the structure sheaf of an affine toric coordinate space. Its path monomials are distinct characters of its split torus. A Hoza weight is a cocharacter, and uniqueness of the minimum produces a one-dimensional initial character with a nonzero coefficient. D-S3 identifies the action and grading; I-10 supplies its flat hypersurface family over a field. Cote~133, scan~019, is only a finite-set analogy here, not a claim that path existence already computes a representation multiplicity.

\paragraph{I-02. Descent through a monomial quotient.}
The monomial map is a map of affine toric schemes whose character-lattice map is the integer matrix $A$. A cocharacter is induced by an integral weight on the full variable torus exactly when its degree vector lies in the image of $A^{\mathsf T}$. Homogeneity of the binomial kernel gives the grading on the image algebra; extending it to the full variable torus additionally requires the integral image-lattice condition in I-02. For example $z\mapsto x^2$ has zero kernel, but weight one on $z$ is not induced by an integer weight on $x$. The diamond is an explicit failure even of kernel homogeneity. This is D-S3, I~4.7.3, applied to the quotient algebra. Cote~133, scan~019 supplies no rule permitting independent path gradings to be identified with the quotient grading.

\paragraph{I-03. A complete coefficient flag.}
Filter the coefficient module by the span of monomials below each cocharacter weight. Basis isolation says that at each growth degree precisely one indexed coefficient survives, and that these surviving degrees are distinct. The dual formulation quantifies over every nonzero quotient functional. This is a statement about filtered locally free modules over a field and their graded pieces, not just the presence of one initial path. D-S3 gives the associated grading, and D-FGA places the open minor conditions on the relevant flag and Grassmann schemes. Cote~133, scan~028's direct-summand question motivates the distinction between a filtration and a chosen splitting.

\paragraph{I-04. The Pl\"ucker weight polytope.}
Take the coefficient row plane in a Grassmannian and its nonzero Pl\"ucker coordinates. Their torus weights are the sums of the corresponding monomial exponents. The vertex-fiber assertion is the symmetric-exchange proof stated in I-04; the probabilistic estimate counts walls crossed in one coordinate. The resulting open cones are normal cones refined by basis-weight inequalities. D-S3 identifies their cocharacters and D-FGA, Expos\'e~221, Proposition~3.8 places the Grassmannian. No finite-field purity theorem chooses an integer point in all these cones automatically.

\paragraph{I-05. Real Gram positivity and algebraic minors.}
With two independent dual coefficient matrices, Cauchy--Binet is an integral identity of sections of a determinant line. On the real locus where one is the conjugate transpose of the other, it becomes a sum of nonnegative squared Pl\"ucker coordinates. The counterexample shows that this determinant line does not retain the whole coefficient-image flag after colliding weights. D-FGA supplies the rank-stratum language, and D-S6 supplies the determinant as top exterior power. Cote~133, scans~019 and~028, respectively mark the distinction between numerical positivity and actual tensor-module information.

\paragraph{I-06. Secant rank of an evaluation functional.}
An evaluation at a point is a rank-one decomposable functional on the bounded-degree polynomial space. The moment matrix of the proposed nonlinear exponent functional has tensor-product rank. Expressing it as a fixed sum of evaluations would place it in a secant locus of bounded rank; the computed determinant excludes that locus below the stated exponential size. This is an algebraic rank statement over $K(t)$, hence belongs to the determinantal geometry in D-FGA. D-S3 also explains why a nonlinear exponent function is not a character-lattice homomorphism. No motive is needed for this counterexample.

\paragraph{I-07. Taking invariants changes the inserted data.}
The invariant functor on complex representations is realized by the Reynolds projector. An independent twirl sends each state insertion to its invariant scalar component. The two-insertion calculation resolves the commutant into the identity and transposition operations and exhibits a negative coefficient. Thus invariants are extracted in the representation category of D-S3/D-S6, with the full tensor operation retained. Cote~133, scan~028, describes why the symmetric-group action on tensor powers matters; scan~019 warns against replacing that action by dimension.

\paragraph{I-08. A finite quotient morphism, conditional on a hitting set.}
The trace-generating functions are regular invariant functions on the affine scheme of matrix tuples. A hitting set separating their coefficient polynomials gives a separating family in the quotient coordinate algebra; the cited invariant-theory theorem makes that algebra finite over the generated subalgebra. This is a finite morphism of affine schemes, not automatically a choice of algebraically independent parameters of minimal cardinality. D-FGA gives the general morphism and parameter-space setting. The unresolved explicitness condition is exactly the one retained in I-08 and GCT~V; the dictionary does not discharge it.

\paragraph{I-09. A tensor contraction computing a Hermitian norm.}
The matrix product computes a regular coefficient map. Tensoring with its conjugate and contracting equal monomial indices gives its squared norm; the torus integral extracts the zero character in that tensor product. D-S3 explains the character extraction. Over an integral algebraic base one must use independent dual variables, not treat complex conjugation as a regular morphism. The metric realization lives over $\mathbb R$ or $\mathbb C$. It is the same coefficient module as before, and access to its presenting matrices remains a white-box assumption.

\paragraph{I-10. Rees modules and specialization.}
The cocharacter gives an increasing filtration on coefficient spans. Its Rees module has the associated graded as its special fiber and the original module on the punctured line. The explicit hypersurface homogenization in I-10 is flat over a field by the torsion-free PID argument in its proof; general ideal degenerations require the entire weight-homogenized ideal. D-S3 identifies the grading and D-FGA identifies the flat family and its Hilbert/Quot points. The associated graded preserves total filtered rank, whereas adding colliding coefficients in a substituted polynomial can reduce their span, as I-05 proves. These are different maps.

\paragraph{I-11. Exterior powers do not remove equations.}
The exterior dimension and graph-of-linear-map count describe the sizes of coordinate and Grassmannian models. The two-term zero determinant is already an equation of the rank-deficient locus. Assigning different path weights to its cancelling terms changes that equation. D-S6 gives exterior powers and D-FGA gives the Grassmannian, but neither turns a signed determinant expansion into unsigned reachability. Cote~133, scan~028's direct-summand margin remains a relevant distinction between a functorial module and its termwise expansion.

\paragraph{I-12. Support is the nonvanishing locus of a regular function.}
A specified coefficient of an integer ROABP is a regular polynomial in its edge parameters. Its nonvanishing open subset is determined by a signed path sum, whereas ordinary reachability detects whether a parameterizing path set is empty. The GapL reduction is a computation of that regular function; the conditional complexity-class inclusion follows only under the specified logarithmic-space reduction. D-FGA locates its zero scheme. Cote~133, scan~019 supplies the nearby conceptual distinction between the cardinality of a finite set and a signed class; it states no complexity theorem.

\paragraph{I-13. Evaluation quadrature and the Segre variety.}
Coefficient tensors of width-one products form the affine cone over a Segre variety. Equality of the quadratic forms on all these tensors, with their complex conjugates, forces the moment-matrix identity proved in I-13. Its rank then bounds the number of evaluation vectors. This restates the entire proof as an assertion about a Segre cone, its dual evaluation vectors, and determinantal rank loci in D-FGA. The Hermitian positivity is not an integral $\lambda$-ring assertion. The much smaller hitting set in the same finding concerns only kernel avoidance, not equality of quadratic forms.

\subsection{Purity, arithmetic sampling, and the limits of numerical shadows}

\paragraph{P-01. A Gram form on a finite space of sections.}
The coefficient support spans a vector space of functions. Their pairwise correlations are entries of a Gram matrix; its spectral lower bound is the elementary estimate in P-01. In geometric applications these functions may be sections or Frobenius traces, but no sheaf is manufactured by naming the Gram matrix. D-ET applies only after an actual trace-function realization exists. The distinction from rank augmentation in cote~133, scan~019, is that all cross terms have been retained and bounded.

\paragraph{P-02. A Kummer-character construction.}
The signs arise from the rank-one quadratic Kummer sheaf pulled back by $t\mapsto t+i$ on the line with its zeros removed. Products over distinct indices remain governed by the explicitly squarefree polynomial, so the curve character estimate controls the off-diagonal trace sums. The altered values at the excluded points contribute the stated elementary error. D-ET, together with the exact Alon--Goldreich--H\aa stad--Peralta locator, carries the proof. This yields the stated sparse multilinear hitting set; it does not create a polynomial bound on arbitrary coefficient supports.

\paragraph{P-03. The kernel of restriction.}
The sign-cube counterexample is a nonzero section whose restriction vanishes on the selected finite subscheme; the common-affine repair is annihilated by the explicit Vandermonde section. These are regular-function identities and survive independently of improved point-counting estimates. They belong to the affine-scheme language of D-FGA. Cote~133, scan~019 is an analogy for information lost by passing to a smaller numerical image, not a source of these new counterexamples.

\paragraph{P-04. Pullback to a curve and its zero divisor.}
For the curve morphism specified by the rational functions $r_i$, the section $f(r_1,\ldots,r_n)$ has polar divisor bounded by the displayed $D$. If it is nonzero, its zeros have total degree at most $\deg D$. The trace formula and curve RH then provide enough rational points. D-ET controls the count; the pullback's nonvanishing is a separate scheme-theoretic condition. This proof uses curves and sheaves actually specified in the statement, so a motivic interpretation is available for their cohomology, not for an unspecified ROABP.

\paragraph{P-05. A curve family can lie in a small polynomial kernel.}
Dimension comparison in the spaces $H^0(C,\mathcal O_C(aD))$ produces an algebraic relation between the two chosen functions. Multiplying the relations for the finite curve family gives its annihilating section and the explicit bivariate ROABP. D-FGA places the image and its ideal; D-ET cannot remove a function from that ideal by increasing the field size. The relation is proved directly with its degree and width bounds and is not a lower bound against every geometric sampler.

\paragraph{P-06. Purity with a complete conductor budget.}
The relevant class is $[\mathcal F_i\otimes\mathcal F_j^\vee]$ in the Grothendieck group of constructible Weil sheaves. The compactly supported trace gives its Gram entry. Geometric irreducibility and inequivalence determine $H_c^2$, affine vanishing removes $H_c^0$, and Weil~II bounds $H_c^1$. SGA~5 supplies the displayed Swan-conductor expression for its dimension. This is exactly D-ET, with every rank and conductor retained. A claim of polynomial explicitness additionally requires that these data and the enumeration of the points be polynomially bounded. Cote~133, scan~028 supplies a tensor-category analogy, but not the missing sheaf realization.

\paragraph{P-07. Sparse cocharacters over the integer lattice.}
The weight family is an explicit finite set in $\operatorname{Hom}(\mathbb G_m,\mathbb G_m^n)=\mathbb Z^n$. The collision polynomial over $\mathbb F_p$ is nonzero by its degree and coefficient bounds, so a seed outside the union of its root sets makes the support characters distinct. D-S3 translates the construction exactly; only the elementary finite-field root bound is used. Thus the result is independent of an additional purity or RH hypothesis.

\paragraph{P-08. Arithmetic covers with a specified nonvanishing encoding.}
The common number field determines a finite \emph{\'etale} cover of $\operatorname{Spec}\mathbb Z$ after inverting its discriminant. Frobenius conjugacy classes are the arithmetic data on that cover. Effective Chebotarev produces a small unramified prime only after the stated common map from primes to evaluations and its nonvanishing property are supplied. D-ET provides the parallel finite-field trace perspective; it is not a proof of that arithmetic encoding. The abelian case uses the stipulated Dirichlet GRH, while the general stated Hecke extension remains conditional on Masullo. No claim about arbitrary automorphic $\mathrm{GL}_n$ GRH is made.

\paragraph{P-Z01 and P-Z02. Realizations of zero statistics.}
These are analytic consequences of the precisely stated RH predicates and supported explicit-formula theorems. Their expression by unitary integrals does not make the zeta zeros the eigenvalues of a constructed geometric Frobenius. D-ET explains the function-field comparison only in the families with an established sheaf and monodromy group. There is no cote~133 scan containing a classical-RH-to-GCT implication; scans~019 and~028 supply representation operations, not that missing implication.

\subsection{The integral controls G-01--G-09}

\paragraph{G-01. Burnside classes, representation classes, and rank.}
The Burnside-ring counterexample is an equivariant finite-set computation. Linearization sends a $G$-set to a permutation module, but the two-subset operation does not become exterior power: the $C_2$ sign in the proof witnesses this failure. In D-S6 the correct relation is $x^2=2\lambda^2(x)+\psi^2(x)$, and only augmentation replaces $\psi^2(x)$ by the rank. This is the precise mathematical correction to cote~133, scan~019, line~53 and the supplied study's equivariant reading.

\paragraph{G-02. The universal terminal variable.}
Specializing one factor to the tensor unit in the universal product identity forces the top elementary symmetric function to remain present. SGA~6, V~2.4, equation~(2.4.2), explicitly includes that terminal variable. The proof in G-02 is its necessity in the universal character ring, so the correction to scan~006, lines~25--26, is mathematical and independent of paleography.

\paragraph{G-03. Equal finite filtration data, unequal representations.}
The two sums of line characters are honest locally free representations of $\mathbb G_{m,\mathbb Z}$, whose difference is the $(D+1)$st power of an augmentation-zero element. D-S6, V~3.10, places the difference in the required filtration step. The associated-bundle construction transports this identity to any scheme carrying the line bundle; it does not force the bundles to remain distinct on every base. Thus finite Chern or $\gamma$ data cannot recover all the character coefficients needed in scans~019 and~028.

\paragraph{G-04. The integral image and its good rank locus.}
The schematic orbit closure is defined over $\mathbb Z$ by a kernel, and its image algebra is torsion-free. In a fixed degree a nonzero maximal minor gives an explicit principal open on which the coefficient matrix is equivalent to a split rank matrix. This is a direct proof of base change for the displayed module, with D-FGA providing the image, Grassmannian and Quot framework. At the excluded primes the flat integral model's fiber may differ from the orbit closure of the reduced point. This is why integral $K_0$ identities from scan~028 do not imply unchanged characteristic-zero multiplicities in all fibers.

\paragraph{G-05. Plane curves and the three-row sector.}
The nonempty complement of the theta divisor in $\operatorname{Pic}^{g-1}(C)$ supplies Beauville's line bundle, hence a linear determinantal equation. Density of smooth curves gives the scheme-theoretic containment needed for every ternary coefficient section. Highest-weight sections of length at most three belong to that coefficient subalgebra, so its determinant ideal has none. D-FGA records the ideal restriction and parameter spaces; D-S3 records the weights. The numerical positive count is not extracted from Chern classes of the curve alone.

\paragraph{G-06. Stabilizer dimensions and four-row equations.}
The polynomial matrix-pencil map is invariant under the explicitly stated subgroup scheme $H$ of $\mathrm{GL}_n^2$. Its generic stabilizer is the scalar subgroup by the concrete matrix calculation. D-S3, VI$_B$~4.1 and its cited EGA theorem justify the dimension behavior, and the resulting proper image has a nonzero ideal. D-FGA and highest-weight restriction give an equation of exactly four rows, using G-05 to exclude shorter weights. This supplies no nonvanishing on a particular padded form until that extra statement is proved.

\paragraph{G-07. A Cartier divisor and a determinant twist.}
The reduced determinantal-quartic locus is a prime hypersurface. Its equation spans the character $\det^D$, so its ideal is the line-character twist of the ambient graded algebra shifted by $D$. The exact sequence in G-07 is the equivariant $K_0$ statement of a principal Cartier divisor. D-S6 gives additivity, D-FGA gives the projective hypersurface and quotient, and D-S3 identifies the character. It produces the shifted plethysm subtraction exactly; the corrected reduced degree is a literature issue, not a new transcription of cote~133.

\paragraph{G-08. A universal morphism of coefficient schemes.}
The two monomial trees and their cofactor determinant give a matrix over the universal integer coefficient ring. Homogenization and block padding therefore commute with any base change. Over $\mathbb C$ the containment of all low-variable restrictions yields an inclusion of highest-weight ideal spaces, hence the multiplicity inequality. This is D-FGA applied to an explicit scheme morphism, with the weights read through D-S3 and the symmetric powers through D-S6. Its complexity bound is part of the construction, not a consequence of representability alone.

\paragraph{G-09. A local padding coordinate and smaller universal matrices.}
Once the nonzero padding form is one coordinate, dehomogenization reduces the number of affine variables by one. The shared exponent chains are again defined over the universal coefficient ring; their exact vertex count gives the stated finite size. Changing an arbitrary padding form to that coordinate is local on its nonzero coefficient charts. D-FGA carries these local morphisms and their closure statement, and the D-S3 highest-weight argument yields the improved row exclusion. The known asymptotic uniform-representation theorem is cited separately; neither the chart construction nor its GCT consequence requires RH.

\subsection{What the manuscript record does and does not establish}
The supplied cote~133 study also reports decisions about the name on scan~008, a characteristic-two pre-$\lambda$ issue on scan~006, and a crossed-out Stiefel--Whitney assertion on scan~009. Those reports are retained as reports of earlier readers, with their own scope. No new visual decision about the name or the deleted text is made in this run. The five attached LaTeX transcriptions contain no original raster pages. The separate transcription table records the forced algebraic readings and the distinction between a textual ambiguity, a mathematical correction, and a claim whose ink cannot be checked here.
